% Figure: capacity-factor comparison across generation technologies.
% Horizontal bars, percent. Representative values current as of 2024.
\begin{figure}[htbp]
\centering
\begin{tikzpicture}
\begin{axis}[
  width=13cm, height=8cm,
  xlabel={capacity factor (\%, current as of 2024)},
  symbolic y coords={solar PV, wind onshore, wind offshore, hydro, coal, gas combined-cycle, geothermal, nuclear},
  ytick=data,
  xmin=0, xmax=100,
  xbar,
  bar width=11pt,
  nodes near coords,
  nodes near coords style={font=\scriptsize},
  tick label style={font=\scriptsize},
  label style={font=\small},
  enlarge y limits=0.12,
]
\addplot[draw=englime!60!black, fill=englime!50] coordinates {
  (18,solar PV)
  (35,wind onshore)
  (45,wind offshore)
  (40,hydro)
  (50,coal)
  (55,gas combined-cycle)
  (80,geothermal)
  (92,nuclear)
};
\end{axis}
\end{tikzpicture}
\caption[Capacity-factor comparison]{Representative capacity factors by
technology (current as of 2024), the ratio of actual annual energy
produced to the energy that would be produced running at rated power
continuously. Intermittent sources sit low; dispatchable thermal sits
middle; nuclear, run as baseload, sits highest. Capacity factor is the
bridge between a plant's rated power and the energy it actually delivers.
Representative figures only.}
\label{fig:vol04ch14:capacity-factor}
\end{figure}
